\begin{annexitem}{annex:p07-is-characterization}{thm:p07-is-characterization}{All four identity-system implications}
\begin{proof}
Call the conditions I, II, III, IV in their displayed order.

I implies II. For \((S,s_0)\), apply induction to
\(D(b,r)=S(b)\), \(d=s_0\); it supplies a pointed map, so
\(\mathsf{ppmap}(R,S)\) is inhabited. Given pointed maps
\((u,\phi)\) and \((v,\psi)\), apply induction to
\(D(b,r)=(u(b,r)=v(b,r))\) with base path
\(\phi\cdot\psi^{-1}\). Obtain
\(h:\prod_{b,r}(u(b,r)=v(b,r))\) and
\(h(a_0,r_0)=\phi\cdot\psi^{-1}\).
Two applications of funext give \(k:u=v\).
Their beta comparisons imply
\(\operatorname{ap}_{w\mapsto w(a_0,r_0)}(k)=h(a_0,r_0)\).
Transport in \(w\mapsto(w(a_0,r_0)=s_0)\) sends \(\phi\) to
\(\operatorname{ap}_{\mathrm{ev}_0}(k)^{-1}\cdot\phi\).
Substitute the displayed base computation. Association, inverse of
a composite and cancellation identify this with
\((\phi\cdot\psi^{-1})^{-1}\cdot\phi=\psi\).
Thus the Sigma-path constructor with first coordinate \(k\)
gives \((u,\phi)=(v,\psi)\). Every two elements are equal, so
the pointed-map type is an inhabited proposition and is contractible.
This step retains equality of the pointing witnesses, not just equality
of their underlying functions.

II implies III. Put \(J_0(b)=(a_0=b)\), pointed by reflexivity.
The canonical map \(e:J_0\to R\) is pointed, since
\(T_R(\mathsf{refl},r_0)\equiv_{\mathrm j}r_0\).
II supplies a pointed map \(d:R\to J_0\) with
\(\delta:d(a_0,r_0)=\mathsf{refl}\).
The composite \(ed\) is pointed by
\(\operatorname{ap}_{e_{a_0}}(\delta)\) and reflexivity transport.
The identity of \(R\) is pointed by reflexivity.
Contractibility of \(\mathsf{ppmap}(R,R)\) identifies these maps,
and application of the first projection followed by evaluation
gives \(e_b(d_b(r))=r\).
For the other composite one must not misapply II to domain \(J_0\).
Based \(J\) gives identity-system induction for
\((J_0,\mathsf{refl})\) itself. The implication I to II, just proved,
therefore makes \(\mathsf{ppmap}(J_0,J_0)\) contractible.
Both \(de\), pointed by \(\delta\), and its identity lie in that
type. Their equality gives \(d_b(e_b(p))=p\).
These two homotopies and adjointification give equivalence of every
specific canonical \(e_b\).

III implies IV. Let
\(U=\sum_b(a_0=b)\), \(T=\sum_bR(b)\).
Define \(F(b,p)=(b,e_b(p))\) and
\(G(b,r)=(b,d_b(r))\), using inverses from III.
Their two homotopies are Sigma paths with first coordinate reflexivity
and second coordinate the inverse homotopy of \(e_b\).
Thus \(F\) is an equivalence by adjointification.
The type \(U\) has center \((a_0,\mathsf{refl})\);
based \(J\) supplies its path to every \((b,p)\), computing to
reflexivity at the center. Apply \(F\) to that contraction at
\(G(z)\), then compose with \(F(G(z))=z\).
Since \(F(a_0,\mathsf{refl})\equiv_{\mathrm j}(a_0,r_0)\),
this is a contraction of \(T\) at the specified point.

IV implies I. Suppose \(T\) has center \(c\) and contraction
\(k_z:c=z\). Let \(z_0=(a_0,r_0)\), and normalize to
\(\kappa_z=k_{z_0}^{-1}\cdot k_z:z_0=z\).
The inverse law gives
\(\kappa_{z_0}=\mathsf{refl}_{z_0}\).
For the uncurried family \(\widehat D:T\to\mathcal V\), define
\(v(z)=T_{\widehat D}(\kappa_z,d)\).
At \(z_0\), application of transport to the preceding loop equality,
then reflexivity transport, gives \(v(z_0)=d\).
Uncurrying and currying supply exactly the requested section and
propositional pointed computation. No truncation elimination or
choice of a merely asserted section is involved.
\end{proof}
\end{annexitem}

\begin{annexitem}{annex:p07-is-two-sided}{thm:p07-is-two-sided}{Two-sided versus local induction and the five criteria}
\begin{proof}
Suppose first that two-sided induction is supplied. Fix \(a_0\),
a local motive \(D:\prod_y(R(a_0,y)\to\mathcal V)\) and
\(d:D(a_0,r_0(a_0))\). Define the global motive
\[
\widetilde D(x,y,r)=
 \prod_{p:x=a_0}
 D\bigl(y,T_{x'\mapsto R(x',y)}(p,r)\bigr).
\]
Its diagonal datum at \(x\) is defined by identity elimination on
\(p:x=a_0\), with \(x\) variable and its right endpoint fixed:
at \(x=a_0,p=\mathsf{refl}\) return \(d\).
This is right-based identity elimination (equivalently ordinary
based \(J\) on the inverse path), not elimination of fixed loops.
It has the stated dependent type because the same elimination
also moves the second endpoint \(y=x\) and the diagonal \(r_0(x)\).
The global eliminator yields \(\widetilde v\).
Define the local section
\(v(y,r)=\widetilde v(a_0,y,r)(\mathsf{refl})\).
Evaluating the global diagonal computation at reflexivity gives
\(v(a_0,r_0(a_0))=d\).

Conversely, for each \(a\), apply its supplied local eliminator to
the motive \(D(a,-,-)\) and datum \(d(a)\).
The resulting sections and base comparisons form dependent functions
of \(a\), so they give the two-sided section and diagonal comparisons.
The eliminators are data indexed by \(a\); no axiom of choice is
used. The local four-way characterization now gives equivalence
with the pointwise canonical-equivalence and total-contractibility
criteria.

We verify also the globally pointed mapping criterion. Currying and
uncurrying give an equivalence between its displayed Sigma type and
\[
\prod_{a:A}
 \mathsf{ppmap}\bigl((R(a,-),r_0(a)),(S(a,-),s_0(a))\bigr).
\]
The forward map evaluates the global function and witness at \(a\);
the backward map collects their first and second projections.
One round trip is judgmental on those projections; in the other,
funext on the functions and on the witness families, then the
Sigma-path constructor, supply the comparison. Thus local identity
systems give a contractible global mapping type: choose the
pointwise centers and use their contractions with funext.

In the reverse direction suppose all global mapping types are
contractible. Take \(S(a,b)=(a=b)\), with diagonal reflexivity, to
obtain a globally pointed \(d:R\to S\).
The canonical transport family \(e:S\to R\) is globally pointed.
Their composite on \(R\) equals its identity by the assumed
contractibility for target \(R\).
For the composite on \(S\), ordinary two-sided \(J\) makes \(S\)
a two-sided identity system; the already proved forward mapping
criterion gives contractibility of its globally pointed self-maps.
Hence this composite too equals the identity. Projection and
evaluation yield both fiberwise inverse homotopies.
Adjointification gives each canonical equivalence, which already
implies the local and global induction criteria.
\end{proof}
\end{annexitem}

\begin{annexitem}{annex:p07-is-equivalence-homotopy-induction}{thm:p07-is-equivalence-homotopy-induction}{The precise axiom-to-induction comparisons}
\begin{proof}
For \(R(X,Y)=(X\simeq Y)\), the canonical transport map from
\(X=Y\), based at the identity equivalence of \(X\), equals
\(\operatorname{idtoeqv}\). To prove the equality eliminate the
variable endpoint universe path; at reflexivity both maps return the
identity equivalence. Univalence asserts that this particular map
is an equivalence. The two-sided characterization therefore supplies
equivalence induction and its propositional diagonal computation.
The universe containing the indexed family and its motive must support
the displayed products; this construction does not lower its level.
Funext is the ambient assumption used by that characterization;
no separate univalence-to-funext proof is inferred here.

For the function type \(X=\prod_aB(a)\), put
\(R(f,g)=\prod_a(f(a)=g(a))\), with diagonal pointwise reflexivity.
The canonical map is \(\mathsf{happly}\): eliminate a function path
\(f=g\), and both maps at reflexivity are
\(a\mapsto\mathsf{refl}_{f(a)}\).
The explicit funext axiom says happly is an equivalence.
Apply the same two-sided characterization to obtain homotopy induction
and its propositional reflexive computation.
Neither conclusion asserts a judgmental reduction at an arbitrary
equivalence or homotopy.
\end{proof}
\end{annexitem}

\begin{annexitem}{annex:p07-flat-pair-calculus}{thm:p07-flat-pair-calculus}{Calibrated dependent paths and every pair-path comparison}
\begin{proof}
Choose \(\mathcal K_s\) by identity elimination on \(s:x=x'\).
At reflexivity, its source is a function identity and its target is
pointwise identity, since horizontal transport computes to the input.
Take the actually assumed happly equivalence with its funext inverse.
This defines both directions at every \(s\), and \(J\) carries the
two beta/eta homotopies to every \(s\).
For constant codomain, normalize constant-family transport by its
J-built comparison; equivalently define \(\mathcal K^0_s\)
by the same \(J\) prescription with happly at reflexivity.
It too has both inverse homotopies. Thus using its inverse followed
by its forward map recovers the supplied component path, rather than
an arbitrary function path with the same type.

For the composition formula, eliminate \(s,r,t,v\), with their
endpoints variable, in that order. The endpoint pairs become one
pair \((x,u)\); \(\gamma_Y\), transport, action on reflexivity and
pair decoding all compute to reflexivity. The two sides of the
formula are then reflexivity. This defines the asserted propositional
comparison for all four inputs. Elimination of \(s,r\) similarly
proves \([s,r]=H(s,u)\cdot V_{x'}(r)\).

For the curried-action formula eliminate \(s\), then \(r\).
The Sigma path is reflexivity, \(L\)'s action is reflexivity,
\(\operatorname{apd}_d(\mathsf{refl})\) is reflexivity,
and \(\mathcal K^0_{\mathsf{refl}}\) is happly; evaluation at \(u\)
therefore returns reflexivity. The last action is also reflexivity,
so their composite is reflexivity. This proves the general formula
by \(J\), with no assumption about setness of the codomain.

Finally for \(\alpha:s=s'\), transport of \(r\) in
\(s\mapsto(T_Y(s,u)=u')\) is
\[
\operatorname{ap}_{t\mapsto T_Y(t,u)}(\alpha)^{-1}\cdot r.
\]
This is the equality-family transport formula, or follows by \(J\)
on \(\alpha\) and the unit paths. A path from that expression to
\(r'\) is precisely the second coordinate needed for a Sigma path
from \((s,r)\) to \((s',r')\). Apply the pair-path decoder to this
Sigma path to obtain a two-path \([s,r]=[s',r']\).
In particular a change of first coordinate alone is insufficient.
\end{proof}
\end{annexitem}

\begin{annexitem}{annex:p07-flat-total-constructors}{thm:p07-flat-total-constructors}{Univalent transport and the two total-space constructors}
\begin{proof}
For the point constructor \(P(\mathsf c(a))\) is judgmentally \(C(a)\),
so \(\mathsf c'(a,y)=(\mathsf c(a),y)\) is well typed.
For \(\mathsf p_b\), pullback transport identifies
\(T_P(\mathsf p_b,y)\) with transport of \(y\) along
\(\operatorname{ap}_P(\mathsf p_b)\) in the identity universe family.
The universe recursor's propositional path computation replaces this
path by \(\operatorname{ua}(D_b)\). The univalence comparison
\(T_{X\mapsto X}(\operatorname{ua}(D_b),y)=D_b(y)\)
then gives the particular path \(\theta_b(y)\).
The Sigma-path decoder applied to
\((\mathsf p_b,\theta_b(y))\) gives exactly the asserted
\(\mathsf p'_{b,y}\), with the displayed endpoints.
The proof does not confuse \(\mathsf p_b\) with a path \(fb=gb\)
in \(A\): its endpoints are in \(W\).
\end{proof}
\end{annexitem}

\begin{annexitem}{annex:p07-flat-total-induction}{thm:p07-flat-total-induction}{The full dependent-function motive and edge lift}
\begin{proof}
Define \(M(w)=\prod_{y:P(w)}Q(w,y)\), where \(Q(w,y)\)
abbreviates \(Q((w,y))\). The point datum is
\(a_a:y\mapsto a_{a,y}\).
For each \(b\), put \(s=\mathsf p_b\), \(y:C(fb)\),
\(t=T_P(s,y)\), \(\theta=\theta_b(y):t=D_b(y)\),
\[
H=H(s,y),\quad V=V_{\mathsf c(gb)}(\theta),\quad
u=T_Q(H,a_{fb,y}),\quad v=a_{gb,t},\quad w=a_{gb,D_b(y)}.
\]
The pair-path calculus gives \(\mathsf p'_{b,y}=H\cdot V\).
Applying transport to this path comparison and then transport
composition gives
\[
\rho:T_Q(\mathsf p'_{b,y},a_{fb,y})=T_Q(V,u).
\]
Hence the supplied edge datum yields
\(\beta=\rho^{-1}\cdot b_{b,y}:T_Q(V,u)=w\).
Dependent action of the fiber section \(z\mapsto a_{gb,z}\),
together with the pullback transport comparison for
\(z\mapsto(\mathsf c(gb),z)\), gives
\(\zeta:T_Q(V,v)=w\).
Here \(V\) is the action of that pair injection on \(\theta\);
this identification is proved by \(J\) on \(\theta\).

For any \(z\) in the source fiber of \(V\), inverse transport supplies
\(\chi_z:T_Q(V^{-1},T_Q(V,z))=z\).
The following composite is therefore explicitly a path \(u=v\):
\[
\delta_{b,y}=
\chi_u^{-1}\cdot
\operatorname{ap}_{T_Q(V^{-1},-)}(\beta\cdot\zeta^{-1})
\cdot\chi_v.
\]
Every factor has its indicated endpoint: the middle factor runs from
the inverse transport of \(T_Q(V,u)\) to that of \(T_Q(V,v)\).
This is the required source-indexed component
\[
T_Q(H(s,y),a_{fb,y})=a_{gb,T_P(s,y)}.
\]
Apply the inverse of \(\mathcal K_s\) to the dependent function
\(y\mapsto\delta_{b,y}\). It supplies
\(T_M(\mathsf p_b,a_{fb})=a_{gb}\), the actual lift for the
coequalizer eliminator.
Induction on \(W\) produces \(d:\prod_wM(w)\) with
\(d(\mathsf c(a))\equiv_{\mathrm j}a_a\).
Set \(v(w,y)=d(w)(y)\); its point computation is the asserted one.
This constructs every required edge lift. It does not assert a
general dependent computation for \(v\) on \(\mathsf p'\), a
stronger statement not needed in the flattening inverse proof.
\end{proof}
\end{annexitem}

\begin{annexitem}{annex:p07-flat-total-recursion}{thm:p07-flat-total-recursion}{A calibrated recursor and its nondependent edge beta}
\begin{proof}
For each \(b,y\), let \(s=\mathsf p_b\),
\(t=T_P(s,y)\) and \(\theta=\theta_b(y):t=D_b(y)\).
Define the explicit component
\[
\delta_{b,y}=b_{b,y}\cdot
  \operatorname{ap}_{z\mapsto a_{gb,z}}(\theta)^{-1}
  :a_{fb,y}=a_{gb,t}.
\]
Use coequalizer induction in \(M(w)=P(w)\to Z\), with point datum
\(a_a\) and lift
\(\mathcal K^0_{\mathsf p_b}{}^{-1}
   (y\mapsto\delta_{b,y})\).
It gives \(d:\prod_w(P(w)\to Z)\), and define \(v(w,y)=d(w)(y)\).
The point computation is judgmental.

For the edge computation, the propositional path computation of
that induction gives
\(\operatorname{apd}_d(\mathsf p_b)
 =\mathcal K^0_{\mathsf p_b}{}^{-1}(\delta_b)\).
Apply \(\mathcal K^0_{\mathsf p_b}\), then its counit, and evaluate
at \(y\). The result is
\[
\mathcal K^0_{\mathsf p_b}
   (\operatorname{apd}_d(\mathsf p_b))(y)=\delta_{b,y}.
\]
The curried-action formula at
\(\mathsf p'_{b,y}=[\mathsf p_b,\theta]\) consequently gives
\[
\operatorname{ap}_v(\mathsf p'_{b,y})
 =\delta_{b,y}\cdot
   \operatorname{ap}_{z\mapsto a_{gb,z}}(\theta)
 =\bigl(b_{b,y}\cdot
   \operatorname{ap}_{a_{gb}}(\theta)^{-1}\bigr)
        \cdot\operatorname{ap}_{a_{gb}}(\theta)
 =b_{b,y}.
\]
The last equality associates to an inverse followed by its path,
cancels that pair, and uses right unit. This is the complete
nondependent path computation, not an assertion that unnamed
transport factors cancel. The calibration through
\(\mathcal K^0\)'s actual counit is essential.
\end{proof}
\end{annexitem}

\begin{annexitem}{annex:p07-flat-equivalence}{thm:p07-flat-equivalence}{Both flattening inverse homotopies and all edge lifts}
\begin{proof}
Let \(T=\sum_wP(w)\).
Recursion on \(\widetilde W\) gives \(h:\widetilde W\to T\)
with
\[
h(\widetilde{\mathsf c}(a,y))
  \equiv_{\mathrm j}\mathsf c'(a,y),\qquad
\operatorname{ap}_h(\widetilde{\mathsf p}_{b,y})
  =\mathsf p'_{b,y}.
\]
Precisely, apply its dependent eliminator to the constant family \(T\),
supplying the constant-family transport comparison followed by
\(\mathsf p'_{b,y}\). Its dependent computation, and
\(\operatorname{apd}=c_{\mathrm{const}}\cdot\operatorname{ap}\),
cancel the same constant-family comparison to give the displayed
ordinary action computation.
The proved total-space recursor gives \(k:T\to\widetilde W\) with
\[
k(\mathsf c'(a,y))\equiv_{\mathrm j}\widetilde{\mathsf c}(a,y),
\qquad \operatorname{ap}_k(\mathsf p'_{b,y})
  =\widetilde{\mathsf p}_{b,y}.
\]

To construct \(kh\sim\mathrm{id}_{\widetilde W}\), use the
dependent motive \(L(z)=(k(h(z))=z)\).
At every point constructor both sides compute to
\(\widetilde{\mathsf c}(a,y)\), so take reflexivity.
For \(e=\widetilde{\mathsf p}_{b,y}\), action on composition and
the two preceding edge computations give
\[
\operatorname{ap}_{kh}(e)
 =\operatorname{ap}_k(\operatorname{ap}_h(e))
 =\operatorname{ap}_k(\mathsf p'_{b,y})
 =e.
\]
Equality-family transport consequently sends the source point
reflexivity to
\(\operatorname{ap}_{kh}(e)^{-1}\cdot\mathsf{refl}\cdot e\).
Substitute the action comparison, use the unit comparisons, then
\(e^{-1}\cdot e=\mathsf{refl}\) at the target point.
This supplies exactly the dependent lift demanded by the
\(\widetilde W\) eliminator. That eliminator gives the first homotopy.

For \(hk\sim\mathrm{id}_T\), use \(N(z)=(h(k(z))=z)\).
Again both functions agree judgmentally at each \(\mathsf c'(a,y)\),
so choose reflexivity there. On \(e=\mathsf p'_{b,y}\),
\[
\operatorname{ap}_{hk}(e)
 =\operatorname{ap}_h(\operatorname{ap}_k(e))
 =\operatorname{ap}_h(\widetilde{\mathsf p}_{b,y})
 =e.
\]
The same equality-family calculation provides
\(T_N(e,\mathsf{refl})=\mathsf{refl}\) at the target point.
Apply the constructed dependent total-space induction theorem with
these point and edge data to obtain the second homotopy.
Only its point computation was required; its unasserted general
dependent edge computation is not being used.
Finally the two maps and these two homotopies give the claimed
equivalence by the routed adjointification theorem.
\end{proof}
\end{annexitem}

\begin{annexitem}{annex:p07-flat-circle-cover}{thm:p07-flat-circle-cover}{The native cover equivalence and every signed contraction case}
\begin{proof}
First account for the unit-indexed presentation, without asserting a
judgmental equality of named HITs. Map its \(W\) to \(S^1\) by
\(\mathsf c(a)\mapsto\mathsf{base}\),
\(\mathsf p_b\mapsto\mathsf{loop}\).
Map the circle back by \(\mathsf{base}\mapsto\mathsf c(\star)\)
and \(\mathsf{loop}\mapsto\mathsf p_\star\).
The composite on the circle agrees at its point and loop by the
recursor computations; equality-family transport and inverse
cancellation give its inverse homotopy by circle induction.
On \(W\), the point comparison
\(\mathsf c(\star)=\mathsf c(a)\) is defined by unit induction on
\(a\), returning reflexivity at \(\star\).
For each \(\mathsf p_b\), unit induction on \(b\) reduces the
needed homotopy lift to
\(\mathsf p_\star^{-1}\cdot\mathsf{refl}\cdot
\mathsf p_\star=\mathsf{refl}\), proved by path laws.
Coequalizer induction gives the other inverse homotopy.
Adjointification gives the presentation equivalence.
For the unit-indexed flattened type, define
\(u(\widetilde{\mathsf c}(a,n))=\mathsf c(n)\) with edge
\(\mathsf p_n\), and define
\(v(\mathsf c(n))=\widetilde{\mathsf c}(\star,n)\) with edge
\(\widetilde{\mathsf p}_{\star,n}\).
On \(R\), \(uv\) has reflexivity point comparisons, and its action
on every \(\mathsf p_n\) equals \(\mathsf p_n\); cancellation
gives the dependent lift.
For \(vu\), define each point comparison
\(\widetilde{\mathsf c}(\star,n)=\widetilde{\mathsf c}(a,n)\)
by unit induction on \(a\). For each edge, unit induction on
\(b\) reduces the lift to the inverse cancellation for
\(\widetilde{\mathsf p}_{\star,n}\), with \(n\) still arbitrary.
The two HIT inductions yield both inverse homotopies, and
adjointification gives equivalence with \(R\).
Thus neither presentation comparison adds judgmental unit eta.

We give the circle specialization directly, so it does not require an
unprovided equality between two recursively defined families.
Let \(T=\sum_x\mathrm{code}(x)\) and
\(S_n=[\mathsf{loop},\theta_n]\), where
\(\theta_n:T_{\mathrm{code}}(\mathsf{loop},n)
 =\mathsf{succ}(n)\) is the already established monodromy.
Recursion on \(R\) gives \(h:R\to T\) with
\(h(\mathsf c(n))\equiv_{\mathrm j}(\mathsf{base},n)\)
and \(\operatorname{ap}_h(\mathsf p_n)=S_n\).
As in the general proof, ordinary action computation is obtained
from dependent constant-family computation by cancelling the
constant-family transport comparison.

For \(k:T\to R\), apply circle induction to
\(M(x)=\mathrm{code}(x)\to R\), with point function \(\mathsf c\).
Its supplied loop lift is
\[
\mathcal K^0_{\mathsf{loop}}{}^{-1}
\left(n\mapsto
 \mathsf p_n\cdot\operatorname{ap}_{\mathsf c}(\theta_n)^{-1}
\right).
\]
The component on the right runs from \(\mathsf c(n)\) to
\(\mathsf c(T_{\mathrm{code}}(\mathsf{loop},n))\), so it has
exactly the required calibrated type. Set \(k(x,u)=d(x)(u)\).
Then \(k(\mathsf{base},n)\equiv_{\mathrm j}\mathsf c(n)\).
The calibrated counit, circle path computation, and curried-action
formula give
\[
\operatorname{ap}_k(S_n)
 =(\mathsf p_n\cdot\operatorname{ap}_{\mathsf c}(\theta_n)^{-1})
       \cdot\operatorname{ap}_{\mathsf c}(\theta_n)
 =\mathsf p_n.
\]

On \(R\), the motive \(k(h(z))=z\) has reflexivity at every
\(\mathsf c(n)\). On \(\mathsf p_n\), action of \(kh\) equals
\(\mathsf p_n\) by the two edge computations above.
Equality-family transport therefore gives the lift by
\(\mathsf p_n^{-1}\cdot\mathsf{refl}\cdot\mathsf p_n
 =\mathsf{refl}\). Induction on \(R\) gives \(kh\sim\mathrm{id}\).
For the other homotopy use \(Q(z)=(h(k(z))=z)\) on \(T\).
Its point datum \(a(n)\) is reflexivity at \((\mathsf{base},n)\).
Action of \(hk\) on \(S_n\) equals \(S_n\), so the same
equality-family calculation supplies
\(\beta_n:T_Q(S_n,a(n))=a(\mathsf{succ}(n))\).
To spell out the total-space induction in its circle form,
write \(t=T_{\mathrm{code}}(\mathsf{loop},n)\),
\(H=H(\mathsf{loop},n)\), \(V=V_{\mathsf{base}}(\theta_n)\).
The pair decomposition \(S_n=H\cdot V\) converts \(\beta_n\)
to \(T_Q(V,T_Q(H,a(n)))=a(\mathsf{succ}(n))\).
Dependent action of \(a\) along \(\theta_n\) gives
\(T_Q(V,a(t))=a(\mathsf{succ}(n))\).
Apply \(T_Q(V^{-1},-)\) and its inverse-transport comparisons
to these two paths, exactly as the explicit \(\delta\) construction
in the total-induction proof. The result is
\(T_Q(H,a(n))=a(t)\).
Applying \(\mathcal K_{\mathsf{loop}}^{-1}\) to these components
provides the circle lift for
\(\prod_{u:\mathrm{code}(x)}Q(x,u)\).
Circle induction gives the section, hence \(hk\sim\mathrm{id}_T\).
Adjointification gives \(T\simeq R\). No circle decoder or either
of its inverse homotopies was used.

For contractibility, define canonical paths
\(\bar q_c:\mathsf c(0)=\mathsf c(E(c))\), with \(c\) in the
signed carrier of the integer presentation:
\[
\begin{aligned}
\bar q_{\mathsf z}&=\mathsf{refl},&
\bar q_{\mathsf{pos}(0)}&=\mathsf{refl}\cdot\mathsf p_0,&
\bar q_{\mathsf{pos}(\mathsf Sk)}
  &=\bar q_{\mathsf{pos}(k)}\cdot\mathsf p_{E(\mathsf{pos}(k))},\\
\bar q_{\mathsf{neg}(0)}
  &=\mathsf{refl}\cdot\mathsf p_{E(\mathsf{neg}(0))}^{-1},&
\bar q_{\mathsf{neg}(\mathsf Sk)}
  &=\bar q_{\mathsf{neg}(k)}
       \cdot\mathsf p_{E(\mathsf{neg}(\mathsf Sk))}^{-1}.
\end{aligned}
\]
Their endpoints are those of the signed successor table; in
particular the negative step uses the edge indexed by the smaller
integer, not by its predecessor's upper endpoint.
For arbitrary \(n\), define
\[
q_n=T_{z\mapsto(\mathsf c(0)=\mathsf c(z))}
          (\nu_n,\bar q_{G(n)}),
\qquad \nu_n:E(G(n))=n.
\]
On each canonical constructor, \(GE\) computes to that constructor.
The path \(\nu_{E(c)}\) is therefore a loop at \(E(c)\).
Integer setness compares it with reflexivity, so
\(q_{E(c)}=\bar q_c\) propositionally. This does not assert
setness of \(R\).

Now prove \(q_n\cdot\mathsf p_n=q_{\mathsf{succ}(n)}\) by
signed induction, using these canonical comparisons.
At zero and at every positive tag, it is the positive defining
clause of \(\bar q\).
At \(\mathsf{neg}(0)\), its left side becomes
\((\mathsf{refl}\cdot\mathsf p_{-1}^{-1})\cdot\mathsf p_{-1}\),
which equals reflexivity by association, inverse cancellation
and units. That is the zero target.
At \(\mathsf{neg}(\mathsf Sk)\), it becomes
\[
(\bar q_{\mathsf{neg}(k)}
       \cdot\mathsf p_{E(\mathsf{neg}(\mathsf Sk))}^{-1})
       \cdot\mathsf p_{E(\mathsf{neg}(\mathsf Sk))}
 =\bar q_{\mathsf{neg}(k)}.
\]
Associate, cancel the inverse pair and use right unit.
These exhaust the successor cases. Signed induction transports
the proved equation family along \(\nu_n\) to arbitrary integers.
Transport in the based-path family \(x\mapsto(\mathsf c(0)=x)\)
along \(\mathsf p_n\) is right concatenation; the just-proved
equation is thus exactly the dependent lift for the \(R\) eliminator.
Its point data \(q_n\) and edge lifts yield a contraction centered
at \(\mathsf c(0)\).
Transport that contraction along the inverse of \(T\simeq R\)
using the two inverse homotopies, as in the proved retract
construction. This makes \(T\) contractible independently of
the encode--decode proof.
\end{proof}
\end{annexitem}

\begin{annexitem}{annex:p07-is-circle-pair-coherence}{thm:p07-is-circle-pair-coherence}{The winding proof's transported second coordinate}
\begin{proof}
For \(n=0\), \(\mathsf{loop}^0\equiv_{\mathrm j}\mathsf{refl}\)
and reflexivity transport computes to \(0\).
Thus \(h_0:0=0\); integer setness gives \(h_0=\mathsf{refl}_0\).
Apply the pair-path decoder to this comparison.
Its computation at two reflexivities gives \(K_0=\mathsf{refl}\).

For arbitrary \(n\), pair composition gives
\[
K_n\cdot S_n=
 [\mathsf{loop}^n\cdot\mathsf{loop},r_n],
\quad
r_n=\gamma_{\mathrm{code}}(\mathsf{loop}^n,\mathsf{loop},0)
 \cdot\operatorname{ap}_{T_{\mathrm{code}}(\mathsf{loop},-)}(h_n)
 \cdot\theta_n.
\]
Its second coordinate has type
\[
r_n:T_{\mathrm{code}}(\mathsf{loop}^n\cdot\mathsf{loop},0)
                         =\mathsf{succ}(n).
\]
The signed power theorem supplies
\(\alpha_n:\mathsf{loop}^n\cdot\mathsf{loop}
                         =\mathsf{loop}^{\mathsf{succ}(n)}\)
(the inverse of its stated successor comparison).
Let
\(A_n=\operatorname{ap}_{s\mapsto T_{\mathrm{code}}(s,0)}(\alpha_n)\).
The paths \(r_n\) and \(A_n\cdot h_{\mathsf{succ}(n)}\)
have the same endpoints in \(\mathbb Z\).
Integer setness gives their equality. Hence, by association and
inverse cancellation,
\[
A_n^{-1}\cdot r_n=h_{\mathsf{succ}(n)}.
\]
This is exactly the transport of the second coordinate along
\(\alpha_n\), by the pair-calculus equality-family formula.
The Sigma-path constructor therefore gives a path between the
two encoded path pairs. Applying the pair-path decoder gives
\[
[\mathsf{loop}^n\cdot\mathsf{loop},r_n]
 =[\mathsf{loop}^{\mathsf{succ}(n)},h_{\mathsf{succ}(n)}]
 =K_{\mathsf{succ}(n)}.
\]
Compose with the original pair composition comparison.
This supplies the two dependencies left to the reader at the
source's Section 8.1 line 771: Sigma composition and the comparison
of the actual winding witnesses after an integer step.
No path proof in the circle itself was declared unique.
\end{proof}
\end{annexitem}

\begin{annexitem}{annex:p07-is-circle-native-induction}{thm:p07-is-circle-native-induction}{The direct transported-data identity-system construction}
\begin{proof}
Let \(T=\sum_x\mathrm{code}(x)\), \(z_0=(\mathsf{base},0)\),
and \(\widehat D:T\to\mathcal V\) be the uncurried motive.
Define the actual base function
\(a(n)=T_{\widehat D}(K_n,d)\).
Transport composition and the proved pair-path coherence give
\[
T_{\widehat D}(S_n,a(n))
 =T_{\widehat D}(K_n\cdot S_n,d)
 =T_{\widehat D}(K_{\mathsf{succ}(n)},d)
 =a(\mathsf{succ}(n)).
\]
The first equality is the inverse of the transport-composition
comparison. The second is application of the transport function
to the previously proved equality of total-space paths.
Call this complete edge path \(\beta_n\).

Here is its dependent circle lift, not just its required type.
Put \(t=T_{\mathrm{code}}(\mathsf{loop},n)\),
\(\theta=\theta_n:t=\mathsf{succ}(n)\),
\(H=H(\mathsf{loop},n)\), \(V=V_{\mathsf{base}}(\theta)\).
Pair decomposition and transport composition convert \(\beta_n\)
to
\[
\beta'_n:
T_{\widehat D}(V,T_{\widehat D}(H,a(n)))
                                  =a(\mathsf{succ}(n)).
\]
Dependent action of \(a\) along \(\theta\), with pullback transport,
gives
\(\zeta_n:T_{\widehat D}(V,a(t))=a(\mathsf{succ}(n))\).
For \(u=T_{\widehat D}(H,a(n))\), \(v=a(t)\), let
\(\chi_z:T_{\widehat D}(V^{-1},T_{\widehat D}(V,z))=z\).
The explicit composite
\[
\chi_u^{-1}\cdot
\operatorname{ap}_{T_{\widehat D}(V^{-1},-)}
      (\beta'_n\cdot\zeta_n^{-1})
\cdot\chi_v
\]
is a path \(T_{\widehat D}(H,a(n))=a(t)\).
Apply \(\mathcal K_{\mathsf{loop}}^{-1}\) to the function
of \(n\) consisting of these paths. It yields the necessary lift
in \(M(x)=\prod_{c:\mathrm{code}(x)}D(x,c)\).
Circle induction with point function \(a\) and this lift produces
the requested section \(v\).

Its point computation gives \(v(\mathsf{base},0)
\equiv_{\mathrm j}T_{\widehat D}(K_0,d)\).
Applying transport to \(K_0=\mathsf{refl}_{z_0}\), then using
reflexivity transport, yields the pointed comparison with \(d\).
Thus \((\mathrm{code},0)\) satisfies identity-system induction.
The general four-way characterization identifies the canonical
map \(p\mapsto T_{\mathrm{code}}(p,0)\) as an equivalence.
This proof used signed winding, not either inverse homotopy of
the earlier decoder, and did not use the cover contraction.
\end{proof}
\end{annexitem}
